\subsection{Four double adjacent hom-sets}
\begin{proposition}[Proved]\label{A7:four-double}
Let $C$ be a finite acyclic category on the totally ordered support
$0<1<2<3<4$. Suppose that each of $C(0,1),C(1,2),C(2,3),C(3,4)$
has exactly two elements. Then the product
$H^2(BC;\mathbb Q)\otimes H^2(BC;\mathbb Q)\longrightarrow H^4(BC;\mathbb Q)$
is zero, regardless of the sizes of the other six hom-sets.
\end{proposition}
\begin{proof}
Write the four adjacent sets as $A=\{a_0,a_1\}$,
$B=\{b_0,b_1\}$, $D=\{c_0,c_1\}$, and $E=\{d_0,d_1\}$.
We use $D,E$ for these sets only, not for categories. Let
\[
 u(a,b)=b\circ a\in C(0,2),\quad
 v(b,c)=c\circ b\in C(1,3),\quad
 w(c,d)=d\circ c\in C(2,4).
\]
Let $\alpha,\beta$ be normalized rational two-cocycles on $NC$.
Every four-simplex has vertices $01234$, and its cup value is
$\alpha(a,b)\beta(c,d)$. If the restriction of $\alpha$ to the
full subcategory on $012$ is zero in cohomology, extend a primitive
one-cochain by zero to $C$ and change $\alpha$ by its coboundary.
The new representative vanishes on all $012$ triangles, proving
$[\alpha][\beta]=0$. The same argument applies if $[\beta]$ restricts
to zero on $234$. We may therefore assume both restrictions nonzero.

The boundary of a two-chain in $C_{012}$ has separate components on
$A$, $B$, and $C(0,2)$. Vanishing of its first two components forces
the chain to be a multiple of
$(a_0-a_1)\otimes(b_0-b_1)$. Consequently
\[
 [u(a_0,b_0)]+[u(a_1,b_1)]
 =[u(a_0,b_1)]+[u(a_1,b_0)]
 \quad\text{in }\mathbb Q[C(0,2)],
\]
and $x=\alpha(a_0,b_0)-\alpha(a_1,b_0)
-\alpha(a_0,b_1)+\alpha(a_1,b_1)$ is nonzero.
Equality of these two multisets says that $u$ is constant on each
column or constant on each row. There are exactly two disjoint cases:
\begin{enumerate}
\item[(U0)] $u(a_0,b)=u(a_1,b)$ for both $b$;
\item[(U1)] $u(a,b_0)=u(a,b_1)$ for both $a$, and the two row values
are distinct.
\end{enumerate}
The analogous argument on $234$ shows that
$y=\beta(c_0,d_0)-\beta(c_0,d_1)
-\beta(c_1,d_0)+\beta(c_1,d_1)\ne0$, and gives two cases:
\begin{enumerate}
\item[(W0)] $w(c,d_0)=w(c,d_1)$ for both $c$;
\item[(W1)] $w(c_0,d)=w(c_1,d)$ for both $d$, with distinct column values.
\end{enumerate}

Let $z$ be any rational four-cycle. The faces omitting $0$ and $4$
land in separate three-chain summands. Their vanishing gives a unique
matrix $(t_{bc})_{b\in B,c\in D}$ such that
\[
 z=\sum_{b,c}t_{bc}
 (a_0-a_1)\otimes b\otimes c\otimes(d_0-d_1).
\]
The remaining three faces imply
\begin{align*}
 &\sum_b t_{bc}([u(a_0,b)]-[u(a_1,b)])=0 &&(c\in D),\\
 &\sum_c t_{bc}([w(c,d_0)]-[w(c,d_1)])=0 &&(b\in B),\\
 &\sum_{b,c}t_{bc}[v(b,c)]=0 &&\text{in }\mathbb Q[C(1,3)].
\end{align*}
Thus U1 forces every column sum of $t$ to vanish, while W1 forces
every row sum to vanish.

Put $\ell_b=\alpha(a_0,b)-\alpha(a_1,b)$ and
$r_c=\beta(c,d_0)-\beta(c,d_1)$, so $x=\ell_{b_0}-\ell_{b_1}$,
$y=r_{c_0}-r_{c_1}$, and
\[
 \langle\alpha\smile\beta,z\rangle=\sum_{b,c}t_{bc}\ell_b r_c.
\]
Define functions on $C(1,3)$ by
$F(s)=\alpha(a_0,s)-\alpha(a_1,s)$ and
$G(s)=\beta(s,d_0)-\beta(s,d_1)$. The cocycle equation
\[
 \gamma(g,h)-\gamma(g\circ f,h)+\gamma(f,h\circ g)-\gamma(f,g)=0
\]
on $0123$ gives $\ell_b=F(v(b,c))$ in case U0, and
$x=F(v(b_0,c))-F(v(b_1,c))$ in case U1. On $1234$ it gives
$r_c=G(v(b,c))$ in case W0, and
$y=G(v(b,c_0))-G(v(b,c_1))$ in case W1.

In case U0/W0 the cup evaluation is
$\sum t_{bc}F(v(b,c))G(v(b,c))=0$, because the coefficients sum to
zero separately over every fibre of $v$.

In case U1/W0 write $t_{b_0c}=q_c$ and $t_{b_1c}=-q_c$.
Applying $F$ to the last boundary equation gives
$x(q_{c_0}+q_{c_1})=0$. Since $x\ne0$, the matrix $t$ is a scalar
multiple of the checkerboard matrix. The same conclusion holds
in case U0/W1 by applying $G$, since $y\ne0$, and it holds directly
in case U1/W1. If that scalar is zero we are done. Otherwise the
last boundary equation says
\[
 [v(b_0,c_0)]+[v(b_1,c_1)]
 =[v(b_0,c_1)]+[v(b_1,c_0)].
\]
Hence $v$ is constant on rows or on columns. If it is constant on
columns, the formulas for $\ell$ (U0) or $x$ (U1) force $x=0$.
If it is constant on rows, the formulas for $r$ (W0) or $y$ (W1)
force $y=0$. Both contradict the nonzero restrictions. Thus the
cup evaluation vanishes on every rational four-cycle. Over a field,
cohomology pairs nondegenerately with homology, proving the claim.
\end{proof}
